\section{Greek and Latin Do Not Preserve the Same Clock}

The Septuagint witness represented by Swete keeps the command to look up and
the sun's approach to setting. At verse 17, however, its printed text has a
flame occurring and does not reproduce the Hebrew's explicit thick darkness.
At verse 18 it prints a spatial ``there'' where the Hebrew has ``on that day.''
The Greek chapter is therefore less temporally insistent at its close.

Augustine's Old Latin citation supplies contrary indirect versional evidence:
his Latin retains ``on that day'' while agreeing with Swete's absence of the
Hebrew darkness phrase at verse 17. It does not by itself identify Augustine's
Greek exemplar or prove variation in Greek transmission. The evidence
authorizes only a claim about Swete's identified printed base text, not about
``the Septuagint'' in every witness.

The Clementine Vulgate preserves the temporal chain: the sun goes down, then
has set, thick darkness appears, and the covenant is made that day. Its
``look up'' and ``brought him outside'' also preserve the spatial movement.
Latin readers therefore receive the difficulty more nearly in the Hebrew form.

